%% evaluation.tex
%%

%% ==================
\section{Evaluation}
\label{sec:Evaluation}
%% ==================

This section presents the experimental results addressing the three research questions: reproduction of the precision collapse (RQ1), effectiveness of the hybrid detector (RQ2), and tunability of the precision/recall trade-off (RQ3). Results are reported as mean $\pm$ standard deviation across 5 independent random seeds (42--46), each with stratified 70/30 train/test splits on 1,200 generated samples.

\subsection{Overall Results}

Table~\ref{tab:results} summarizes the performance of all detector variants, bound to \texttt{results/results\_summary.csv} (per-seed rows in \texttt{results\_per\_seed.csv}). The key findings are: (1) ML detector achieves perfect performance with rich context but degrades substantially in code-only setting, (2) hybrid detector restores perfect precision in code-only setting but at significant recall cost, and (3) rule-based detector establishes the recall floor that the hybrid cannot exceed.

\begin{table}[htp]
\centering
\caption{Results across 5 seeds (mean $\pm$ std)}
\label{tab:results}
\begin{tabular}{lccc}
\hline
Model & Precision & Recall & F1 \\
\hline
ML Detector (rich context) & 1.000 $\pm$ 0.000 & 1.000 $\pm$ 0.000 & 1.000 $\pm$ 0.000 \\
ML Detector (code-only, baseline gap) & 0.768 $\pm$ 0.045 & 0.708 $\pm$ 0.064 & 0.733 $\pm$ 0.019 \\
Rule-Based Only (reference) & 1.000 $\pm$ 0.000 & 0.483 $\pm$ 0.028 & 0.651 $\pm$ 0.025 \\
Hybrid Detector (improved, code-only) & 1.000 $\pm$ 0.000 & 0.483 $\pm$ 0.028 & 0.651 $\pm$ 0.025 \\
\hline
\end{tabular}
\end{table}

%% ===============================
\subsection{RQ1: Reproducing the Precision Collapse}
\label{sec:evalOne}
%% ===============================

\textbf{Research Question:} Can we reproduce the precision collapse finding from War et al.~\cite{War25} on a controlled synthetic benchmark where the ``easy'' vs. ``hard'' distinction is explicit by construction?

\textbf{Answer:} Yes. The ML detector trained and evaluated on rich-context (with-comments) data achieves perfect performance: precision 1.000, recall 1.000, F1 1.000, with zero variance across all 5 seeds. This is expected given how the benchmark is constructed: comment language directly names the risk (``WARNING: hardcoded fallback password''), providing an unambiguous signal to a text classifier.

When those same comments are stripped (code-only setting), performance degrades substantially:
\begin{itemize}
  \item Precision drops from 1.000 to 0.768 ($\pm$ 0.045), a \textbf{23.2\% relative decrease}
  \item Recall drops from 1.000 to 0.708 ($\pm$ 0.064), a \textbf{29.2\% relative decrease}
  \item F1 drops from 1.000 to 0.733 ($\pm$ 0.019), a \textbf{26.7\% relative decrease}
\end{itemize}

This reproduces the qualitative direction of War et al.'s own ablation study: removing natural-language context causes the ML classifier to over-flag (precision collapse) while simultaneously missing some genuine misconfigurations (recall degradation). The fact that precision and recall both degrade—rather than a precision/recall trade-off where one improves as the other worsens—indicates that the model is simply less confident and less accurate in the absence of comment signals.

\textbf{Per-Seed Stability:} The standard deviations are small relative to the means (e.g., precision: 0.045 / 0.768 = 5.9\% coefficient of variation), indicating that the phenomenon is robust across different random samples and not an artifact of a single lucky or unlucky seed.

\textbf{Interpretation:} The degradation is not a failure of the ML detector per se—it is performing as well as can be expected given the information available. The hard patterns are genuinely indistinguishable from code alone, and the ML detector has no mechanism to recognize this and abstain. It must predict something, and without the comment signal, its predictions are unreliable.

%% ===============================
\subsection{RQ2: The Hybrid Fix -- and Its Real Cost}
\label{sec:evalTwo}
%% ===============================

\textbf{Research Question:} Does a hybrid rule-based + ML detector restore precision in the code-only setting, and if so, at what cost to recall?

\textbf{Answer:} The hybrid detector fully restores precision to 1.000 ($\pm$ 0.000) in the code-only setting, matching the rich-context ideal. However, this comes at a substantial cost: recall is 0.483 ($\pm$ 0.028), identical to the rule-based-only reference and \textbf{31.7\% lower} than the plain code-only ML classifier's recall of 0.708.

\subsubsection{Why Precision is Restored}

The hybrid detector's perfect precision is achieved through two complementary mechanisms:

\begin{enumerate}
  \item \textbf{Rule Layer Handles Easy Patterns:} For snippets containing an easy misconfiguration (e.g., \texttt{mode = "0777"}), the rule layer fires immediately, guaranteeing a correct positive prediction. Because the rule patterns are constructed to match only genuine misconfigurations, these predictions have zero false positives.
  
  \item \textbf{Strict Threshold Filters Unreliable ML Predictions:} For snippets the rule layer does not recognize, the ML classifier's predicted probability must exceed 0.8 to trigger a positive prediction. This strict threshold filters out the low-confidence guesses that drove false positives in the plain code-only ML detector. Empirically, nearly all easy-pattern misconfigurations that the rule layer misses (there are none by construction in this benchmark, but in principle could exist) and a small fraction of hard-pattern misconfigurations exceed this threshold.
\end{enumerate}

\subsubsection{Why Recall is Constrained}

The hybrid detector's recall is structurally bounded by the rule layer's coverage. Breaking down recall by pattern type:

\begin{itemize}
  \item \textbf{Easy Patterns (50\% of samples):} The rule layer detects all easy misconfigurations (recall = 1.000). The ML fallback is never invoked because the rule fires first.
  
  \item \textbf{Hard Patterns (50\% of samples):} The rule layer detects none of these (recall = 0.000). The ML classifier, trained on code-only data, has no signal to work with—the code line is identical for safe and misconfigured hard patterns. As a result, its predicted probabilities for hard-pattern misconfigurations are effectively random and rarely exceed 0.8.
\end{itemize}

Overall recall is approximately $0.5 \times 1.0 + 0.5 \times 0.0 = 0.5$, with the measured value of 0.483 reflecting the empirical 50/50 hard/easy split plus noise from the random filler lines.

\textbf{Comparison with Rule-Only:} The hybrid detector's recall (0.483) is statistically indistinguishable from the rule-only detector's recall (0.483), confirming that the ML fallback contributes negligible recall gain at the default threshold of 0.8. The hybrid detector's advantage is that it offers a \emph{mechanism} to trade recall for precision via threshold tuning, whereas the rule-only detector has no such knob.

\textbf{Comparison with Code-Only ML:} The plain code-only ML detector achieves higher recall (0.708) by making low-confidence predictions on hard patterns, but this comes at the cost of precision (0.768). The hybrid detector makes the opposite trade-off: it sacrifices recall on hard patterns to achieve perfect precision on easy patterns.

%% ===============================
\subsection{RQ3: The Trade-off Is Tunable, Not Fixed}
\label{sec:evalThree}
%% ===============================

\textbf{Research Question:} Is the precision/recall trade-off tunable via the ML confidence threshold, or is it a fixed operating point?

\textbf{Answer:} The trade-off is tunable. Sweeping the hybrid detector's ML confidence threshold from 0.5 (equivalent to the plain ML detector) to 0.9 (extremely strict) traces a clear Pareto frontier:

\begin{center}
\begin{tabular}{lccc}
\hline
Threshold & Precision & Recall & F1 \\
\hline
0.5 (= plain ML) & 0.768 & 0.708 & 0.733 \\
0.6 & 0.928 & 0.528 & 0.672 \\
0.7 & 0.998 & 0.485 & 0.652 \\
0.8 (default) & 1.000 & 0.483 & 0.651 \\
0.9 & 1.000 & 0.483 & 0.651 \\
\hline
\end{tabular}
\end{center}

\subsubsection{Observations}

\begin{enumerate}
  \item \textbf{Precision Rises Monotonically:} As the threshold increases, precision improves from 0.768 (threshold 0.5) to 1.000 (threshold 0.8), then plateaus. This confirms that stricter thresholds successfully filter out unreliable ML predictions.
  
  \item \textbf{Recall Falls to Rule-Layer Floor:} Recall decreases from 0.708 (threshold 0.5, where the hybrid behaves like the plain ML detector plus rule boosts) down to 0.483 (threshold 0.8 and above, where the ML fallback contributes almost nothing). It does not drop below 0.483 because the rule layer's contribution is threshold-independent—it always catches easy patterns.
  
  \item \textbf{F1 Optimum at Threshold 0.5:} Maximum F1 (0.733) occurs at threshold 0.5, but this operating point offers no advantage over the plain code-only ML detector. The hybrid's value proposition is precision recovery at higher thresholds, not F1 maximization.
  
  \item \textbf{Precision Ceiling at Threshold 0.8:} Beyond 0.8, precision remains at 1.000 and recall remains at 0.483, indicating that the remaining false positives below threshold 0.8 have already been eliminated. Further threshold increases provide no benefit.
\end{enumerate}

\subsubsection{Operational Implications}

An operator deploying this system must choose a threshold based on their risk tolerance:

\begin{itemize}
  \item \textbf{High-Precision Regime (threshold $\geq$ 0.8):} Zero false positives, suitable for automated blocking or alerting where false alarms are costly. Recall is limited to easy patterns only. Recommended for production environments where precision is paramount.
  
  \item \textbf{Balanced Regime (threshold 0.6--0.7):} Precision 0.93--1.00, recall 0.49--0.53. Some false positives remain but are rare. Suitable for semi-automated review workflows where human analysts can triage flagged snippets.
  
  \item \textbf{High-Recall Regime (threshold 0.5):} Equivalent to plain ML detector. Precision 0.768, recall 0.708. Suitable for initial screening or research contexts where missing misconfigurations is more costly than false positives.
\end{itemize}

\textbf{No Free Lunch:} Critically, there is no threshold that achieves both high precision \emph{and} high recall. The fundamental constraint is that hard patterns are genuinely ambiguous in code-only form. A hybrid detector can choose where on the Pareto frontier to operate, but it cannot escape the frontier itself.

%% ===============================
\subsection{Breakdown by Pattern Type}
\label{sec:pattern-breakdown}
%% ===============================

To further validate that the recall limitation is due to hard patterns (not a failure of the rule layer on easy patterns), we examined the per-pattern-type performance on a single representative seed (seed 42):

\begin{center}
\begin{tabular}{lcccc}
\hline
Pattern Type & Rule Recall & ML Recall & Hybrid Recall & Ground Truth \\
\hline
Easy (all 4 combined) & 1.000 & 1.000 & 1.000 & Code-distinguishable \\
Hard (both combined) & 0.000 & 0.410 & 0.000 & Comment-dependent \\
\hline
\end{tabular}
\end{center}

This confirms:
\begin{itemize}
  \item The rule layer has perfect recall on easy patterns (1.000), as expected by construction.
  \item The rule layer has zero recall on hard patterns (0.000), also as expected—there is no fixed literal to match.
  \item The code-only ML classifier achieves 0.410 recall on hard patterns despite having no true signal, indicating it is making educated guesses based on filler line correlations or other spurious features. These predictions are unreliable (precision would be low if evaluated on hard patterns alone).
  \item The hybrid detector at threshold 0.8 has zero recall on hard patterns, confirming that the ML fallback does not contribute at this threshold.
\end{itemize}

%% ===============================
\subsection{Discussion}
\label{sec:discussion}
%% ===============================

\subsubsection{Summary of Findings}

This project successfully reproduces War et al.'s precision collapse finding on a controlled synthetic benchmark (RQ1), builds a hybrid detector that restores perfect precision in the code-only setting (RQ2), and demonstrates that the precision/recall trade-off is tunable via confidence thresholding (RQ3). The hybrid detector is effective for its intended use case: ensuring zero false positives on enumerable misconfigurations while gracefully degrading on context-dependent cases.

\subsubsection{Honest Limitations}

The hybrid detector's recall limitation (0.483 vs. the code-only ML detector's 0.708) is not a bug—it is a structural consequence of the benchmark design and the decision to prioritize precision. The hard patterns are genuinely indistinguishable from code alone, and no amount of rule/ML combination can recover information that was never in the code to begin with.

This is exactly the point: the baseline paper's own conclusion is that semantic context (comments, task descriptions) is \emph{load-bearing}, not merely helpful. If a misconfiguration's safety depends on information that exists only in a comment, then comment quality in IaC repositories is itself a security control, not just a readability nicety. Automated detection tools can only go so far.

\subsubsection{External Validity}

The synthetic benchmark used here has a known 50/50 hard/easy ratio by construction. Real-world IaC repositories likely have different ratios, and may have entirely different ``hard'' pattern categories (e.g., misconfigurations that depend on cross-file context, architectural knowledge, or deployment environment). The hybrid detector's effectiveness in practice depends on:

\begin{enumerate}
  \item The prevalence of enumerable (rule-detectable) vs. context-dependent misconfigurations in the target codebase
  \item The quality and coverage of the rule set (this project uses 4 simple patterns; production systems like Checkov use hundreds)
  \item The availability and quality of natural-language comments in real IaC scripts
\end{enumerate}

Future work should evaluate the hybrid approach on real-world datasets (Ansible Galaxy, Puppet Forge, Terraform Registry) with ground-truth annotations to measure its practical effectiveness. The synthetic benchmark establishes proof-of-concept; real-world validation would establish production readiness.

\subsubsection{Comparison with Baseline Paper}

War et al.~\cite{War25} identify the precision collapse but do not propose or evaluate a fix. This project contributes:
\begin{itemize}
  \item A concrete solution (hybrid detector) for the subset of misconfigurations that are rule-representable
  \item An explicit characterization of the solution's limitations (hard patterns remain unsolvable)
  \item A tunable trade-off curve, enabling operators to choose their preferred operating point
\end{itemize}

The baseline paper's recommendation is to use rich semantic context whenever possible—this project accepts that recommendation and addresses the fallback case where comments are sparse or absent.

\subsubsection{Implications for Practitioners}

For practitioners developing or using IaC security tools:
\begin{enumerate}
  \item \textbf{Comment Quality Matters:} If your IaC scripts contain security-critical configuration decisions that are not self-evident from the code (e.g., why a particular secret reference is safe), document them explicitly. Automated tools cannot infer intent.
  
  \item \textbf{Hybrid Detectors Are Not Silver Bullets:} Combining rules and ML can improve precision, but cannot recover recall on genuinely ambiguous cases. Set realistic expectations about detection coverage.
  
  \item \textbf{Threshold Tuning Is Essential:} The default threshold (0.8 in this project) represents one point on a trade-off curve. Operators should tune the threshold based on their false positive vs. false negative cost ratio.
  
  \item \textbf{Rule Coverage Is Key:} The hybrid detector's recall is bounded by the rule layer's coverage. Investing in comprehensive rule sets (as tools like Checkov do) directly improves recall.
\end{enumerate}
